\chapter{12.11-2 等边三角形的判定}



\section*{12.11.2 等边三角形的判定}

\subsection*{一、选择题（共 8 小题）}

\questionlist

\item 已知 $a, b, c$ 是 $\triangle ABC$ 的三边，且 $a^2 + b^2 + c^2 = ab + ac + bc$，则 $\triangle ABC$ 是（\quad）

  \optionlist
  \item 等腰三角形
  \item 直角三角形
  \item 等边三角形
  \item 等腰直角三角形
  \end{enumerate}

\item 下列三角形：

\begin{enumerate}[label=(\arabic*)]
\item 有两个角等于 $60^\circ$；
\item 有一个角等于 $60^\circ$ 的等腰三角形；
\item 三个外角（每个顶点处各取一个外角）都相等的三角形；
\item 一腰上的中线也是这条腰上的高的等腰三角形。
\end{enumerate}

其中是等边三角形的有（\quad）

\optionlist
\item 1,2,3
\item 1,2,4
\item 1,3
\item 1,2,3,4
\end{enumerate}

\item
  \begin{questionwithimage}{12-11-2-3.jpg}
  如图，$E$ 是等边 $\triangle ABC$ 中 $AC$ 边上的点，$\angle 1 = \angle 2$，$BE = CD$，则 $\triangle ADE$ 的形状是（\quad）
  \optionlist
  \item 等腰三角形
  \item 等边三角形
  \item 不等边三角形
  \item 不能确定形状
  \end{enumerate}
  \end{questionwithimage}

\item 下列推理错误的是（\quad）

\optionlist
\item 在 $\triangle ABC$ 中，$\because \angle A = \angle B = \angle C$，$\therefore \triangle ABC$ 为等边三角形
\item 在 $\triangle ABC$ 中，$AB = AC$，且 $\angle B = \angle C$，$\triangle ABC$ 为等边三角形
\item 在 $\triangle ABC$ 中，$\angle A = 60^\circ$，$\angle B = 60^\circ$，$\triangle ABC$ 为等边三角形
\item 在 $\triangle ABC$ 中，$AB = AC$，$\angle B = 60^\circ$，$\triangle ABC$ 为等边三角形
\end{enumerate}

\item 三角形中有两条中线分别平分它的两个内角，则这个三角形是（\quad）

  \optionlist
  \item 直角三角形
  \item 等腰三角形
  \item 等边三角形
  \item 等腰直角三角形
  \end{enumerate}

\item
  \begin{questionwithimage}{12-11-2-6.png}
  如图，等边 $\triangle DEF$ 的顶点分别在等边 $\triangle ABC$ 的各边上，且 $DE \perp BC$ 于 $E$，若 $AB = 1$，则 $DB$ 的长为（\quad）
  \optionlist
  \item $\frac{1}{2}$
  \item $\frac{1}{3}$
  \item $\frac{2}{3}$
  \item $\frac{3}{4}$
  \end{enumerate}
  \end{questionwithimage}

\item
  \begin{questionwithimage}{12-11-2-7.png}
  如图，在等边三角形 $ABC$ 中，$BC = 2$，$D$ 是 $AB$ 的中点，过点 $D$ 作 $DF \perp AC$ 于点 $F$，过点 $F$ 作 $EF \perp BC$ 于点 $E$，则 $BE$ 的长为（\quad）
  \optionlist
  \item 1
  \item $\frac{3}{2}$
  \item $\frac{5}{4}$
  \item $\frac{4}{3}$
  \end{enumerate}
  \end{questionwithimage}

\item 如图，$\angle MON = 30^\circ$。点 $A_1, A_2, A_3, \ldots$ 在射线 $ON$ 上，点 $B_1, B_2, B_3, \ldots$ 在射线 $OM$ 上。$\triangle A_1B_1A_2$、$\triangle A_2B_2A_3$、$\triangle A_3B_3A_4$、$\ldots$ 均为等边三角形。若 $OA_1 = 1$，则 $\triangle A_7B_7A_8$ 的边长为（\quad）
  \begin{questionwithimage}{12-11-2-8.png}
  \optionlist
  \item 6
  \item 12
  \item 32
  \item 64
  \end{enumerate}
  \end{questionwithimage}

\end{enumerate}

\section*{二、填空题（共 6 小题）}

\blanklist[9]
\item
  若正三角形的边长为 $2\sqrt{5}\,cm$，则这个正三角形的高是\underline{\hspace{2cm}}

\item
  在平面直角坐标系中，等边$\triangle ABC$的顶点$A$（-6，0)，$B$（2，0)，则顶点$C$的坐标为\underline{\hspace{2cm}}

\item
  $\triangle ABC$ 是等边三角形，点$D$是$BC$边上任意一点， $DE \perp AB$ 于点$E$， $DF \perp AC$ 于点$F$. 若 $BC=$ 4，则 $BE+CF=$\underline{\hspace{2cm}}

\item
  已知a、b、c是 $\triangle ABC$ 的三边的长，且满足 $a^{2}+2b^{2}+c^{2}-2b(a+c)=0$，则此三角形的形状为\underline{\hspace{2cm}}

\item
  如图， $\triangle ABC$ 是一个边长为1的等边三角形， $BB_{1}$ 是 $\triangle ABC$ 的高，$B_{1}B_{2}$ 是 $\triangle ABB_{1}$ 的高 $B_{2}B_{3}$ 是 $\triangle AB_{1}B_{2}$ 的高， $\ldots , \quad B_{n-1}B_{n}$ 是 $\triangle AB_{n-2}B_{n-1}$ 的高，$B_{n-1}B_{n}$ 的长是\underline{\hspace{2cm}}

  \centerquestionimage{0.4\textwidth}{12-11-3-1.jpg}

\item
  在 $\triangle ABC$ 中， $\angle C=90^{\circ}$ $AB=2\,cm$ $BC=1\,cm$ .以点$C$为顶点作一个等边三角形，使其他两个顶点在 $\triangle ABC$ 的边上，则这个等边三角形的面积为\underline{\hspace{2cm}}

\end{enumerate}

\section*{三、解答题（共 1 小题）}

\solutionlist[15]
\item
  在 $\mathrm{Rt}\triangle ABC$ 中， $\angle ACB=90^{\circ}$，BD是 $\triangle ABC$ 的角平分线.
(1) 如图1，若 $\scriptstyle AD=BD$，求 $\angle A$ 的度数；
(2)如图2，在（1)的条件下，作 $DE \perp AB$ 于E，连接 $EC$. 求证： $\triangle EBC$ 是等边三角形.

  \begin{questionimagegroup}
  \questionimageitem{0.32\textwidth}{12-11-3-2.jpg}{图1}
  \questionimagegap
  \questionimageitem{0.32\textwidth}{12-11-3-3.jpg}{图2}
  \end{questionimagegroup}

\item
  如图，在 $\triangle ABC$ 中，$\angle A=60^{\circ}$，$BE \perp AC$，垂足为 $E$，$CF \perp AB$，垂足为 $F$，点 $D$ 是 $BC$ 的中点，$BE$，$CF$ 交于点 $M$。

  (1) 如果 $AB=AC$，求证：$\triangle DEF$ 是等边三角形；

  (2) 如果 $AB\ne AC$，试猜想 $\triangle DEF$ 是不是等边三角形？如果 $\triangle DEF$ 是等边三角形，请加以证明；如果 $\triangle DEF$ 不是等边三角形，请说明理由；

  (3) 如果 $CM=4$，$FM=5$，求 $BE$ 的长度。

  \begin{questionimagegroup}
    \questionimageitem[width=\linewidth,height=0.22\textheight,keepaspectratio]{0.32\textwidth}{12-11-2-16-1.png}{}
  \questionimagegap[0.04\textwidth]
    \questionimageitem[width=\linewidth,height=0.22\textheight,keepaspectratio]{0.32\textwidth}{12-11-2-16-2.png}{}
  \end{questionimagegroup}

\vspace{5cm}

\item
  如图所示，$\triangle OAB$ 是边长为 $2+\sqrt{3}$ 的等边三角形，其中 $O$ 是坐标原点，顶点 $A$ 在 $x$ 轴的正方向上。将 $\triangle OAB$ 折叠，使点 $B$ 落在边 $OA$ 上，记为 $B'$，折痕为 $EF$。

  (1) 设 $OB'$ 的长为 $x$，$\triangle OB'E$ 的周长为 $c$，用含有 $x$ 的式子表示 $c$；

  (2) 当 $B'E \parallel y$ 轴时，求点 $B'$ 和点 $E$ 的坐标；

  (3) 当 $B'$ 在 $OA$ 上运动但不与 $O$、$A$ 重合时，能否使 $\triangle EB'F$ 成为直角三角形？若能，请求出点 $B'$ 的坐标；若不能，请说明理由。

\centerquestionimage[width=0.35\textwidth]{0.35\textwidth}{12-11-2-17-1.png}

\vspace{5cm}

\item 已知 $\triangle ABC$ 是等腰三角形，$AB = AC$，$D$ 为边 $BC$ 上任意一点，$DE \perp AB$ 于 $E$，$DF \perp AC$ 于 $F$，且 $E, F$ 分别在边 $AB, AC$ 上。

(1) 如图 a，当 $\triangle ABC$ 是等边三角形时，证明：$AE + AF = \frac{3}{2}BC$。

(2) 如图 b，若 $\triangle ABC$ 中，$\angle BAC = 120^\circ$，探究线段 $AE, AF, AB$ 之间的数量关系，并对你的猜想加以证明。

(3) 如图 c，若 $\triangle ABC$ 中，$AB = 10, BC = 16, EF = 6$，利用你对 (1), (2) 两题的解题思路计算出线段 $CD$（$BD > CD$）的长。

\centerquestionimage[width=0.60\textwidth]{0.60\textwidth}{12-11-2-18.png}


\end{enumerate}
